% Options for packages loaded elsewhere
\PassOptionsToPackage{unicode}{hyperref}
\PassOptionsToPackage{hyphens}{url}
\documentclass[
  11pt,
  a4paper,
]{article}
\usepackage{xcolor}
\usepackage[margin=18mm]{geometry}
\usepackage{amsmath,amssymb}
\setcounter{secnumdepth}{-\maxdimen} % remove section numbering
\usepackage{iftex}
\ifPDFTeX
  \usepackage[T1]{fontenc}
  \usepackage[utf8]{inputenc}
  \usepackage{textcomp} % provide euro and other symbols
\else % if luatex or xetex
  \usepackage{unicode-math} % this also loads fontspec
  \defaultfontfeatures{Scale=MatchLowercase}
  \defaultfontfeatures[\rmfamily]{Ligatures=TeX,Scale=1}
\fi
\ifPDFTeX\else
  % xetex/luatex font selection
    \setmainfont[]{Noto Serif}
    \setsansfont[]{Noto Sans}
    \setmonofont[]{Noto Sans Mono}
  \setmathfont[]{Noto Sans Math}
\fi
% Use upquote if available, for straight quotes in verbatim environments
\IfFileExists{upquote.sty}{\usepackage{upquote}}{}
\IfFileExists{microtype.sty}{% use microtype if available
  \usepackage[]{microtype}
  \UseMicrotypeSet[protrusion]{basicmath} % disable protrusion for tt fonts
}{}
\makeatletter
\@ifundefined{KOMAClassName}{% if non-KOMA class
  \IfFileExists{parskip.sty}{%
    \usepackage{parskip}
  }{% else
    \setlength{\parindent}{0pt}
    \setlength{\parskip}{6pt plus 2pt minus 1pt}}
}{% if KOMA class
  \KOMAoptions{parskip=half}}
\makeatother
\ifLuaTeX
\usepackage[bidi=basic]{babel}
\else
\usepackage[bidi=default]{babel}
\fi
\babelprovide[main,import]{russian}
\ifPDFTeX
\else
\babelfont{rm}[]{Noto Serif}
\fi
% get rid of language-specific shorthands (see #6817):
\let\LanguageShortHands\languageshorthands
\def\languageshorthands#1{}
\setlength{\emergencystretch}{3em} % prevent overfull lines
\providecommand{\tightlist}{%
  \setlength{\itemsep}{0pt}\setlength{\parskip}{0pt}}
\usepackage{bookmark}
\IfFileExists{xurl.sty}{\usepackage{xurl}}{} % add URL line breaks if available
\urlstyle{same}
\hypersetup{
  pdftitle={Билеты по дискретным случайным процессам},
  pdflang={ru-RU},
  hidelinks,
  pdfcreator={LaTeX via pandoc}}

\title{Билеты по дискретным случайным процессам}
\author{}
\date{}

\begin{document}
\maketitle

\section{\texorpdfstring{Билет 04. Борелевские множества, борелевские
\(\sigma\)-алгебры, пополнение меры и борелевские меры
Лебега}{Билет 04. Борелевские множества, борелевские \textbackslash sigma-алгебры, пополнение меры и борелевские меры Лебега}}\label{ux431ux438ux43bux435ux442-04.-ux431ux43eux440ux435ux43bux435ux432ux441ux43aux438ux435-ux43cux43dux43eux436ux435ux441ux442ux432ux430-ux431ux43eux440ux435ux43bux435ux432ux441ux43aux438ux435-sigma-ux430ux43bux433ux435ux431ux440ux44b-ux43fux43eux43fux43eux43bux43dux435ux43dux438ux435-ux43cux435ux440ux44b-ux438-ux431ux43eux440ux435ux43bux435ux432ux441ux43aux438ux435-ux43cux435ux440ux44b-ux43bux435ux431ux435ux433ux430}

Источники: Ширяев 1, гл. II; лекции ДСП.

\subsection{1. Короткий
ответ}\label{ux43aux43eux440ux43eux442ux43aux438ux439-ux43eux442ux432ux435ux442}

Если \(X\) - топологическое пространство, то борелевская
\(\sigma\)-алгебра на \(X\) - это минимальная \(\sigma\)-алгебра,
содержащая все открытые множества:

\[
\mathcal B(X)=\sigma\{G\subset X:\ G\text{ открыто}\}.
\]

Ее элементы называются борелевскими множествами. В \(\mathbb R^d\)
борелевскую \(\sigma\)-алгебру можно порождать открытыми шарами,
открытыми множествами, замкнутыми множествами, прямоугольниками или
полуинтервалами:

\[
\mathcal B(\mathbb R^d)
=
\sigma\left\{\prod_{j=1}^d(a_j,b_j):a_j<b_j\right\}.
\]

Борелевская мера - это мера, заданная на \(\mathcal B(X)\). Борелевская
мера Лебега на \(\mathbb R^d\) - это мера на
\(\mathcal B(\mathbb R^d)\), которая на прямоугольниках равна обычному
объему:

\[
\lambda_d^B\left(\prod_{j=1}^d[a_j,b_j)\right)
=
\prod_{j=1}^d(b_j-a_j).
\]

Мера называется полной, если любое подмножество множества меры ноль
измеримо:

\[
N\subset Z,\quad Z\in\mathcal F,\quad \mu(Z)=0
\quad\Longrightarrow\quad
N\in\mathcal F.
\]

Пополнение меры добавляет к \(\mathcal F\) все подмножества нулевых
множеств. Полная мера Лебега - это пополнение борелевской меры Лебега.
Поэтому каждое борелевское множество лебегово измеримо, но не каждое
лебегово измеримое множество является борелевским.

\subsection{2. Полный
ответ}\label{ux43fux43eux43bux43dux44bux439-ux43eux442ux432ux435ux442}

\subsubsection{2.1. Введение: Соотношение борелевских и лебеговых
множеств}\label{ux432ux432ux435ux434ux435ux43dux438ux435-ux441ux43eux43eux442ux43dux43eux448ux435ux43dux438ux435-ux431ux43eux440ux435ux43bux435ux432ux441ux43aux438ux445-ux438-ux43bux435ux431ux435ux433ux43eux432ux44bux445-ux43cux43dux43eux436ux435ux441ux442ux432}

В предыдущем билете мера Лебега строилась как продолжение
геометрического объема. Теперь нужно различить две близкие, но не
одинаковые системы множеств: борелевские множества и лебегово измеримые
множества.

\subsubsection{\texorpdfstring{2.2. Определение борелевской
\(\sigma\)-алгебры}{2.2. Определение борелевской \textbackslash sigma-алгебры}}\label{ux43eux43fux440ux435ux434ux435ux43bux435ux43dux438ux435-ux431ux43eux440ux435ux43bux435ux432ux441ux43aux43eux439-sigma-ux430ux43bux433ux435ux431ux440ux44b}

Борелевская \(\sigma\)-алгебра появляется из топологии. Если на
множестве \(X\) задана топология, то известны открытые множества. Но для
меры и вероятности открытых множеств недостаточно: нужна
\(\sigma\)-алгебра, чтобы брать дополнения, счетные объединения и
счетные пересечения. Поэтому берут минимальную \(\sigma\)-алгебру,
содержащую все открытые множества:

\[
\mathcal B(X)=\sigma(\text{открытые множества}).
\]

Такую \(\sigma\)-алгебру называют борелевской, а ее элементы -
борелевскими множествами.

На прямой и в \(\mathbb R^d\) это основной класс множеств для
распределений. Например, если \(X\) - случайная величина, то ее
распределение является вероятностной мерой на \(\mathcal B(\mathbb R)\).
Если \(\xi\) - случайный вектор в \(\mathbb R^d\), то его распределение
является мерой на \(\mathcal B(\mathbb R^d)\).

\subsubsection{\texorpdfstring{2.3. Порождающие классы в
\(\mathbb R^d\)}{2.3. Порождающие классы в \textbackslash mathbb R\^{}d}}\label{ux43fux43eux440ux43eux436ux434ux430ux44eux449ux438ux435-ux43aux43bux430ux441ux441ux44b-ux432-mathbb-rd}

В \(\mathbb R^d\) борелевская \(\sigma\)-алгебра может быть порождена
разными удобными классами:

\[
\mathcal B(\mathbb R^d)
=
\sigma(\text{открытые шары})
=
\sigma(\text{замкнутые множества})
=
\sigma(\text{полуоткрытые прямоугольники}).
\]

Это важно практически: функцию распределения удобно задавать через лучи
\((-\infty,x]\), меру Лебега - через прямоугольники, а топологические
рассуждения - через открытые множества. Все эти подходы ведут к одной и
той же борелевской \(\sigma\)-алгебре.

\subsubsection{2.4. Борелевская мера и проблема
полноты}\label{ux431ux43eux440ux435ux43bux435ux432ux441ux43aux430ux44f-ux43cux435ux440ux430-ux438-ux43fux440ux43eux431ux43bux435ux43cux430-ux43fux43eux43bux43dux43eux442ux44b}

Борелевская мера - это просто мера на \(\mathcal B(X)\). Например,
борелевская мера Лебега \(\lambda_d^B\) на \(\mathbb R^d\) задается как
продолжение объема прямоугольников на \(\mathcal B(\mathbb R^d)\):

\[
\lambda_d^B\left(\prod_{j=1}^d[a_j,b_j)\right)
=
\prod_{j=1}^d(b_j-a_j).
\]

Но борелевская мера Лебега не полна. Это значит, что может существовать
борелевское множество \(Z\) меры ноль и подмножество \(N\subset Z\),
которое не является борелевским. Тогда \(N\) должно иметь меру ноль с
геометрической точки зрения, но оно не лежит в области определения
борелевской меры.

\subsubsection{2.5. Процедура пополнения
меры}\label{ux43fux440ux43eux446ux435ux434ux443ux440ux430-ux43fux43eux43fux43eux43bux43dux435ux43dux438ux44f-ux43cux435ux440ux44b}

Чтобы исправить это, меру пополняют. Пусть есть пространство с мерой
\((\Omega,\mathcal F,\mu)\). Полное пополнение имеет \(\sigma\)-алгебру

\[
\overline{\mathcal F}^{\,\mu}
=
\{A\cup N:\ A\in\mathcal F,\ N\subset Z,\ Z\in\mathcal F,\ \mu(Z)=0\}.
\]

Идея: берем старые измеримые множества \(A\) и разрешаем добавлять к ним
произвольные подмножества старых нулевых множеств. Мера на таком
множестве определяется по формуле

\[
\overline{\mu}(A\cup N)=\mu(A).
\]

Это корректно: добавление или удаление подмножеств нулевых множеств не
должно менять меру.

\subsubsection{2.6. Мера Лебега и иерархия
измеримости}\label{ux43cux435ux440ux430-ux43bux435ux431ux435ux433ux430-ux438-ux438ux435ux440ux430ux440ux445ux438ux44f-ux438ux437ux43cux435ux440ux438ux43cux43eux441ux442ux438}

Полная мера Лебега \(\lambda_d\) - это пополнение борелевской меры
Лебега \(\lambda_d^B\). Поэтому есть строгая цепочка:

\[
\mathcal B(\mathbb R^d)
\subset
\mathcal L(\mathbb R^d)
\subset
2^{\mathbb R^d},
\]

где \(\mathcal L(\mathbb R^d)\) - \(\sigma\)-алгебра лебегово измеримых
множеств. Первое включение строго: существуют лебегово измеримые
множества, которые не являются борелевскими. Второе включение тоже
строго: существуют неизмеримые по Лебегу множества, например множество
Витали.

\subsection{3. Основные
определения}\label{ux43eux441ux43dux43eux432ux43dux44bux435-ux43eux43fux440ux435ux434ux435ux43bux435ux43dux438ux44f}

\subsubsection{Топологическое
пространство}\label{ux442ux43eux43fux43eux43bux43eux433ux438ux447ux435ux441ux43aux43eux435-ux43fux440ux43eux441ux442ux440ux430ux43dux441ux442ux432ux43e}

Пара \((X,\tau)\), где \(\tau\) - семейство открытых подмножеств \(X\),
замкнутое относительно произвольных объединений и конечных пересечений.

\subsubsection{\texorpdfstring{Борелевская
\(\sigma\)-алгебра}{Борелевская \textbackslash sigma-алгебра}}\label{ux431ux43eux440ux435ux43bux435ux432ux441ux43aux430ux44f-sigma-ux430ux43bux433ux435ux431ux440ux430}

\[
\mathcal B(X)=\sigma(\tau)
=
\sigma\{G\subset X:\ G\text{ открыто}\}.
\]

Это минимальная \(\sigma\)-алгебра, содержащая все открытые множества.

\subsubsection{Борелевское
множество}\label{ux431ux43eux440ux435ux43bux435ux432ux441ux43aux43eux435-ux43cux43dux43eux436ux435ux441ux442ux432ux43e}

Элемент \(\mathcal B(X)\).

\subsubsection{Борелевская
мера}\label{ux431ux43eux440ux435ux43bux435ux432ux441ux43aux430ux44f-ux43cux435ux440ux430}

Мера \(\mu\) на измеримом пространстве \((X,\mathcal B(X))\).

\subsubsection{Борелевская вероятностная
мера}\label{ux431ux43eux440ux435ux43bux435ux432ux441ux43aux430ux44f-ux432ux435ux440ux43eux44fux442ux43dux43eux441ux442ux43dux430ux44f-ux43cux435ux440ux430}

Борелевская мера \(\mu\) такая, что

\[
\mu(X)=1.
\]

\subsubsection{Борелевская мера
Лебега}\label{ux431ux43eux440ux435ux43bux435ux432ux441ux43aux430ux44f-ux43cux435ux440ux430-ux43bux435ux431ux435ux433ux430}

Мера \(\lambda_d^B\) на \(\mathcal B(\mathbb R^d)\), совпадающая с
объемом на полуоткрытых прямоугольниках:

\[
\lambda_d^B\left(\prod_{j=1}^d[a_j,b_j)\right)
=
\prod_{j=1}^d(b_j-a_j).
\]

\subsubsection{Полная
мера}\label{ux43fux43eux43bux43dux430ux44f-ux43cux435ux440ux430}

Мера \(\mu\) на \((\Omega,\mathcal F)\) называется полной, если

\[
N\subset Z,\quad Z\in\mathcal F,\quad \mu(Z)=0
\quad\Longrightarrow\quad
N\in\mathcal F.
\]

\subsubsection{\texorpdfstring{Пополнение \(\sigma\)-алгебры по
мере}{Пополнение \textbackslash sigma-алгебры по мере}}\label{ux43fux43eux43fux43eux43bux43dux435ux43dux438ux435-sigma-ux430ux43bux433ux435ux431ux440ux44b-ux43fux43e-ux43cux435ux440ux435}

\[
\overline{\mathcal F}^{\,\mu}
=
\{A\cup N:\ A\in\mathcal F,\ N\subset Z,\ Z\in\mathcal F,\ \mu(Z)=0\}.
\]

\subsubsection{Пополненная
мера}\label{ux43fux43eux43fux43eux43bux43dux435ux43dux43dux430ux44f-ux43cux435ux440ux430}

Если \(E=A\cup N\in\overline{\mathcal F}^{\,\mu}\), где
\(A\in\mathcal F\) и \(N\) лежит в измеримом нулевом множестве, то

\[
\overline{\mu}(E)=\mu(A).
\]

\subsubsection{\texorpdfstring{Лебегова
\(\sigma\)-алгебра}{Лебегова \textbackslash sigma-алгебра}}\label{ux43bux435ux431ux435ux433ux43eux432ux430-sigma-ux430ux43bux433ux435ux431ux440ux430}

Пополнение борелевской \(\sigma\)-алгебры \(\mathcal B(\mathbb R^d)\)
относительно борелевской меры Лебега:

\[
\mathcal L(\mathbb R^d)
=
\overline{\mathcal B(\mathbb R^d)}^{\,\lambda_d^B}.
\]

\subsubsection{Мера
Лебега}\label{ux43cux435ux440ux430-ux43bux435ux431ux435ux433ux430}

Пополнение борелевской меры Лебега. Обозначается \(\lambda_d\).

\subsection{4. Основные теоремы и
свойства}\label{ux43eux441ux43dux43eux432ux43dux44bux435-ux442ux435ux43eux440ux435ux43cux44b-ux438-ux441ux432ux43eux439ux441ux442ux432ux430}

\subsubsection{\texorpdfstring{Свойство 1. Борелевская
\(\sigma\)-алгебра
существует}{Свойство 1. Борелевская \textbackslash sigma-алгебра существует}}\label{ux441ux432ux43eux439ux441ux442ux432ux43e-1.-ux431ux43eux440ux435ux43bux435ux432ux441ux43aux430ux44f-sigma-ux430ux43bux433ux435ux431ux440ux430-ux441ux443ux449ux435ux441ux442ux432ux443ux435ux442}

Для любого топологического пространства \(X\) существует минимальная
\(\sigma\)-алгебра, содержащая все открытые множества.

\subsubsection{Доказательство}\label{ux434ux43eux43aux430ux437ux430ux442ux435ux43bux44cux441ux442ux432ux43e}

Рассмотрим все \(\sigma\)-алгебры на \(X\), содержащие все открытые
множества. Такой класс не пуст, потому что \(2^X\) - \(\sigma\)-алгебра
и содержит все подмножества \(X\). Пересечение любого семейства
\(\sigma\)-алгебр снова является \(\sigma\)-алгеброй. Поэтому

\[
\mathcal B(X)=
\bigcap\{\mathcal F:\ \mathcal F\text{ - }\sigma\text{-алгебра},\ \tau\subset\mathcal F\}
\]

является \(\sigma\)-алгеброй, содержит все открытые множества и
содержится в любой другой \(\sigma\)-алгебре с этим свойством. Значит
она минимальна. \(\square\)

\subsubsection{Свойство 2. Замкнутые множества
борелевские}\label{ux441ux432ux43eux439ux441ux442ux432ux43e-2.-ux437ux430ux43cux43aux43dux443ux442ux44bux435-ux43cux43dux43eux436ux435ux441ux442ux432ux430-ux431ux43eux440ux435ux43bux435ux432ux441ux43aux438ux435}

Если \(F\subset X\) замкнуто, то \(F\in\mathcal B(X)\).

\subsubsection{Доказательство}\label{ux434ux43eux43aux430ux437ux430ux442ux435ux43bux44cux441ux442ux432ux43e-1}

Если \(F\) замкнуто, то \(\overline{F}=X\setminus F\) открыто. Открытое
множество лежит в \(\mathcal B(X)\), а \(\sigma\)-алгебра замкнута
относительно дополнений. Значит \(F\in\mathcal B(X)\). \(\square\)

\subsubsection{\texorpdfstring{Свойство 3. В пространстве со счетной
базой борелевская \(\sigma\)-алгебра порождена
базой}{Свойство 3. В пространстве со счетной базой борелевская \textbackslash sigma-алгебра порождена базой}}\label{ux441ux432ux43eux439ux441ux442ux432ux43e-3.-ux432-ux43fux440ux43eux441ux442ux440ux430ux43dux441ux442ux432ux435-ux441ux43e-ux441ux447ux435ux442ux43dux43eux439-ux431ux430ux437ux43eux439-ux431ux43eux440ux435ux43bux435ux432ux441ux43aux430ux44f-sigma-ux430ux43bux433ux435ux431ux440ux430-ux43fux43eux440ux43eux436ux434ux435ux43dux430-ux431ux430ux437ux43eux439}

Пусть \(\{U_n\}_{n\ge1}\) - счетная база топологии \(X\). Тогда

\[
\mathcal B(X)=\sigma(U_1,U_2,\ldots).
\]

\subsubsection{Доказательство}\label{ux434ux43eux43aux430ux437ux430ux442ux435ux43bux44cux441ux442ux432ux43e-2}

Так как все \(U_n\) открыты, то

\[
\sigma(U_1,U_2,\ldots)\subset\mathcal B(X).
\]

Обратно, любое открытое множество \(G\) представляется как объединение
базовых элементов:

\[
G=\bigcup_{n:\ U_n\subset G}U_n.
\]

Так как база счетная, это счетное объединение. Значит
\(G\in\sigma(U_1,U_2,\ldots)\) для любого открытого \(G\).
Следовательно, \(\mathcal B(X)\subset\sigma(U_1,U_2,\ldots)\). Получаем
равенство. \(\square\)

\subsubsection{\texorpdfstring{Свойство 4. Порождающие классы для
\(\mathcal B(\mathbb R^d)\)}{Свойство 4. Порождающие классы для \textbackslash mathcal B(\textbackslash mathbb R\^{}d)}}\label{ux441ux432ux43eux439ux441ux442ux432ux43e-4.-ux43fux43eux440ux43eux436ux434ux430ux44eux449ux438ux435-ux43aux43bux430ux441ux441ux44b-ux434ux43bux44f-mathcal-bmathbb-rd}

В \(\mathbb R^d\):

\[
\mathcal B(\mathbb R^d)
=
\sigma(\text{открытые шары})
=
\sigma(\text{замкнутые множества})
=
\sigma(\text{полуоткрытые прямоугольники}).
\]

Можно также брать только шары или прямоугольники с рациональными
координатами.

\subsubsection{Доказательство-идея}\label{ux434ux43eux43aux430ux437ux430ux442ux435ux43bux44cux441ux442ux432ux43e-ux438ux434ux435ux44f}

Открытые шары порождают топологию. Замкнутые множества порождают ту же
\(\sigma\)-алгебру, потому что дополнение замкнутого множества открыто,
а дополнение открытого замкнуто. Прямоугольники с рациональными
координатами образуют счетную базу обычной топологии в \(\mathbb R^d\).
По свойству 3 они порождают борелевскую \(\sigma\)-алгебру.

\subsubsection{\texorpdfstring{Свойство 5. Борелевская
\(\sigma\)-алгебра на прямой порождается
лучами}{Свойство 5. Борелевская \textbackslash sigma-алгебра на прямой порождается лучами}}\label{ux441ux432ux43eux439ux441ux442ux432ux43e-5.-ux431ux43eux440ux435ux43bux435ux432ux441ux43aux430ux44f-sigma-ux430ux43bux433ux435ux431ux440ux430-ux43dux430-ux43fux440ux44fux43cux43eux439-ux43fux43eux440ux43eux436ux434ux430ux435ux442ux441ux44f-ux43bux443ux447ux430ux43cux438}

На \(\mathbb R\):

\[
\mathcal B(\mathbb R)=\sigma\{(-\infty,x]:x\in\mathbb R\}.
\]

\subsubsection{Доказательство}\label{ux434ux43eux43aux430ux437ux430ux442ux435ux43bux44cux441ux442ux432ux43e-3}

Каждый луч \((-\infty,x]\) замкнут, значит борелевский. Поэтому

\[
\sigma\{(-\infty,x]\}\subset\mathcal B(\mathbb R).
\]

Обратно, интервалы вида \((a,b]\) выражаются через лучи:

\[
(a,b]=(-\infty,b]\setminus(-\infty,a].
\]

Открытый интервал можно представить как счетное объединение:

\[
(a,b)=\bigcup_{n=1}^{\infty}\left(a,b-\frac1n\right]
\]

для достаточно больших \(n\); точнее можно брать все \(n\), для которых
\(b-\frac1n>a\). Значит открытые интервалы лежат в
\(\sigma\{(-\infty,x]\}\). Любое открытое множество на прямой является
счетным объединением открытых интервалов с рациональными концами,
поэтому оно также лежит в этой \(\sigma\)-алгебре. Следовательно,

\[
\mathcal B(\mathbb R)\subset\sigma\{(-\infty,x]\}.
\]

Получаем равенство. \(\square\)

\subsubsection{Теорема 6. Пополнение
меры}\label{ux442ux435ux43eux440ux435ux43cux430-6.-ux43fux43eux43fux43eux43bux43dux435ux43dux438ux435-ux43cux435ux440ux44b}

Пусть \((\Omega,\mathcal F,\mu)\) - пространство с мерой. Тогда

\[
\overline{\mathcal F}^{\,\mu}
=
\{A\cup N:\ A\in\mathcal F,\ N\subset Z,\ Z\in\mathcal F,\ \mu(Z)=0\}
\]

является \(\sigma\)-алгеброй, а формула

\[
\overline{\mu}(A\cup N)=\mu(A)
\]

задает на ней полную меру. Кроме того, \(\overline{\mu}\) продолжает
\(\mu\).

\subsubsection{Доказательство-идея}\label{ux434ux43eux43aux430ux437ux430ux442ux435ux43bux44cux441ux442ux432ux43e-ux438ux434ux435ux44f-1}

Если множество отличается от измеримого множества только на подмножество
нулевого множества, то его можно считать измеримым после пополнения.
Удобная эквивалентная форма:

\[
E\in\overline{\mathcal F}^{\,\mu}
\quad\Longleftrightarrow\quad
\exists A\in\mathcal F:\ E\triangle A\subset Z,\ Z\in\mathcal F,\ \mu(Z)=0,
\]

где \(E\triangle A\) - симметрическая разность. Из этой формы легко
проверить замкнутость относительно дополнений и счетных объединений.
Если

\[
E_n\triangle A_n\subset Z_n,\qquad \mu(Z_n)=0,
\]

то

\[
\left(\bigcup_n E_n\right)\triangle\left(\bigcup_n A_n\right)
\subset
\bigcup_n Z_n,
\]

а \(\bigcup_n Z_n\) имеет меру \(0\) по счетной субаддитивности. Значит
счетное объединение снова лежит в пополнении.

Корректность \(\overline{\mu}\) следует из того, что два возможных
измеримых представителя отличаются только на множестве меры \(0\),
поэтому имеют одинаковую меру. Полнота очевидна из построения: все
подмножества старых и новых нулевых множеств добавлены.

\subsubsection{Свойство 7. Борелевская мера Лебега не
полна}\label{ux441ux432ux43eux439ux441ux442ux432ux43e-7.-ux431ux43eux440ux435ux43bux435ux432ux441ux43aux430ux44f-ux43cux435ux440ux430-ux43bux435ux431ux435ux433ux430-ux43dux435-ux43fux43eux43bux43dux430}

Борелевская мера Лебега \(\lambda_d^B\) на \(\mathcal B(\mathbb R^d)\)
не является полной.

\subsubsection{Обоснование}\label{ux43eux431ux43eux441ux43dux43eux432ux430ux43dux438ux435}

Достаточно рассмотреть \(d=1\). Канторово множество \(C\) замкнуто,
значит борелевское, и имеет меру Лебега \(0\). Все подмножества \(C\)
имеют внешнюю меру \(0\) и входят в лебегово пополнение. Но подмножеств
у \(C\) больше, чем борелевских множеств на прямой: борелевских множеств
имеет мощность континуума, а подмножеств канторова множества -
\(2^{\mathfrak c}\). Поэтому существуют подмножества \(C\), которые не
являются борелевскими. Они лежат в лебеговом пополнении, но не лежат в
\(\mathcal B(\mathbb R)\). Значит борелевская мера Лебега не полна.

\subsubsection{\texorpdfstring{Свойство 8. Лебегова \(\sigma\)-алгебра
строго шире
борелевской}{Свойство 8. Лебегова \textbackslash sigma-алгебра строго шире борелевской}}\label{ux441ux432ux43eux439ux441ux442ux432ux43e-8.-ux43bux435ux431ux435ux433ux43eux432ux430-sigma-ux430ux43bux433ux435ux431ux440ux430-ux441ux442ux440ux43eux433ux43e-ux448ux438ux440ux435-ux431ux43eux440ux435ux43bux435ux432ux441ux43aux43eux439}

На \(\mathbb R^d\):

\[
\mathcal B(\mathbb R^d)\subsetneq\mathcal L(\mathbb R^d).
\]

При этом

\[
\mathcal L(\mathbb R^d)\ne 2^{\mathbb R^d}.
\]

Первое строгое включение следует из предыдущего свойства. Второе
означает существование неизмеримых по Лебегу множеств, например множеств
Витали.

\subsection{5. Примеры}\label{ux43fux440ux438ux43cux435ux440ux44b}

\subsubsection{Пример 1. Интервалы
борелевские}\label{ux43fux440ux438ux43cux435ux440-1.-ux438ux43dux442ux435ux440ux432ux430ux43bux44b-ux431ux43eux440ux435ux43bux435ux432ux441ux43aux438ux435}

На прямой открытые интервалы борелевские по определению, замкнутые
интервалы борелевские как замкнутые множества, а полуинтервалы
борелевские, потому что

\[
(a,b]=(-\infty,b]\setminus(-\infty,a].
\]

\subsubsection{Пример 2. Счетные множества
борелевские}\label{ux43fux440ux438ux43cux435ux440-2.-ux441ux447ux435ux442ux43dux44bux435-ux43cux43dux43eux436ux435ux441ux442ux432ux430-ux431ux43eux440ux435ux43bux435ux432ux441ux43aux438ux435}

Каждая точка \(\{x\}\) замкнута, значит борелевская. Если
\(A=\{x_1,x_2,\ldots\}\) счетно, то

\[
A=\bigcup_{n=1}^{\infty}\{x_n\},
\]

поэтому \(A\) борелевско. В частности, \(\mathbb Q\) и
\(\mathbb Q\cap[0,1]\) борелевские.

\subsubsection{Пример 3. Счетные множества имеют меру Лебега
ноль}\label{ux43fux440ux438ux43cux435ux440-3.-ux441ux447ux435ux442ux43dux44bux435-ux43cux43dux43eux436ux435ux441ux442ux432ux430-ux438ux43cux435ux44eux442-ux43cux435ux440ux443-ux43bux435ux431ux435ux433ux430-ux43dux43eux43bux44c}

Так как \(\lambda(\{x\})=0\), то

\[
\lambda(\mathbb Q\cap[0,1])
\le
\sum_{q\in\mathbb Q\cap[0,1]}\lambda(\{q\})
=0.
\]

Значит

\[
\lambda(\mathbb Q\cap[0,1])=0.
\]

\subsubsection{Пример 4. Подмножества нулевых множеств после
пополнения}\label{ux43fux440ux438ux43cux435ux440-4.-ux43fux43eux434ux43cux43dux43eux436ux435ux441ux442ux432ux430-ux43dux443ux43bux435ux432ux44bux445-ux43cux43dux43eux436ux435ux441ux442ux432-ux43fux43eux441ux43bux435-ux43fux43eux43fux43eux43bux43dux435ux43dux438ux44f}

Пусть \(C\) - канторово множество. Оно замкнуто, значит борелевское, и

\[
\lambda(C)=0.
\]

Любое \(N\subset C\) входит в лебегово пополнение и имеет меру \(0\). Но
не каждое такое \(N\) борелевское.

\subsubsection{Пример 5. Неизмеримое
множество}\label{ux43fux440ux438ux43cux435ux440-5.-ux43dux435ux438ux437ux43cux435ux440ux438ux43cux43eux435-ux43cux43dux43eux436ux435ux441ux442ux432ux43e}

Множество Витали \(V\subset[0,1]\) не является лебегово измеримым.
Значит даже после пополнения меры Лебега область определения не
становится равной \(2^{\mathbb R}\).

\subsection{6. Доказательства и
обоснования}\label{ux434ux43eux43aux430ux437ux430ux442ux435ux43bux44cux441ux442ux432ux430-ux438-ux43eux431ux43eux441ux43dux43eux432ux430ux43dux438ux44f}

\textbf{Почему борелевская \(\sigma\)-алгебра зависит от топологии.}

Определение

\[
\mathcal B(X)=\sigma(\text{открытые множества})
\]

использует открытые множества. Если на одном и том же \(X\) взять другую
топологию, то изменится класс открытых множеств, а значит может
измениться и борелевская \(\sigma\)-алгебра.

\textbf{Почему в \(\mathbb R^d\) достаточно рациональных
прямоугольников.}

Обычная топология в \(\mathbb R^d\) имеет счетную базу из
прямоугольников с рациональными координатами. Любое открытое множество
является объединением всех таких базовых прямоугольников, которые в нем
лежат. Так как база счетна, это счетное объединение. Поэтому
\(\sigma\)-алгебра, порожденная рациональными прямоугольниками, уже
содержит все открытые множества.

\textbf{Почему пополнение не меняет старую меру.}

Если \(A\in\mathcal F\), то можно представить \(A=A\cup\varnothing\),
где \(\varnothing\) лежит в любом нулевом множестве. Тогда

\[
\overline{\mu}(A)=\mu(A).
\]

Поэтому пополнение является именно продолжением меры, а не новой мерой с
другими значениями на старых множествах.

\textbf{Почему добавляются только подмножества нулевых множеств.}

Если \(N\subset Z\) и \(\mu(Z)=0\), то по монотонности любая разумная
мера должна давать

\[
0\le\mu(N)\le\mu(Z)=0.
\]

Значит единственно возможное значение для \(N\) - ноль. Поэтому такие
множества можно добавить без противоречия. Для произвольных подмножеств
положительных множеств такой аргумент не работает.

\textbf{Почему не все лебегово измеримые множества борелевские.}

Пополнение добавляет все подмножества борелевских нулевых множеств. У
канторова множества мера \(0\), но подмножеств у него больше, чем всех
борелевских множеств на прямой. Поэтому среди этих подмножеств есть
неборелевские, но после пополнения они становятся лебегово измеримыми.

\subsection{7. Что написать на
доске}\label{ux447ux442ux43e-ux43dux430ux43fux438ux441ux430ux442ux44c-ux43dux430-ux434ux43eux441ux43aux435}

\begin{enumerate}
\def\labelenumi{\arabic{enumi}.}
\tightlist
\item
  Борелевская \(\sigma\)-алгебра:
\end{enumerate}

\[
\mathcal B(X)=\sigma\{G\subset X:\ G\text{ открыто}\}.
\]

\begin{enumerate}
\def\labelenumi{\arabic{enumi}.}
\setcounter{enumi}{1}
\tightlist
\item
  В \(\mathbb R^d\):
\end{enumerate}

\[
\mathcal B(\mathbb R^d)
=
\sigma(\text{открытые шары})
=
\sigma(\text{полуоткрытые прямоугольники}).
\]

\begin{enumerate}
\def\labelenumi{\arabic{enumi}.}
\setcounter{enumi}{2}
\tightlist
\item
  Борелевская мера Лебега:
\end{enumerate}

\[
\lambda_d^B\left(\prod_{j=1}^d[a_j,b_j)\right)
=
\prod_{j=1}^d(b_j-a_j).
\]

\begin{enumerate}
\def\labelenumi{\arabic{enumi}.}
\setcounter{enumi}{3}
\tightlist
\item
  Полная мера:
\end{enumerate}

\[
N\subset Z,\quad Z\in\mathcal F,\quad \mu(Z)=0
\Rightarrow
N\in\mathcal F.
\]

\begin{enumerate}
\def\labelenumi{\arabic{enumi}.}
\setcounter{enumi}{4}
\tightlist
\item
  Пополнение:
\end{enumerate}

\[
\overline{\mathcal F}^{\,\mu}
=
\{A\cup N:\ A\in\mathcal F,\ N\subset Z,\ Z\in\mathcal F,\ \mu(Z)=0\}.
\]

\begin{enumerate}
\def\labelenumi{\arabic{enumi}.}
\setcounter{enumi}{5}
\tightlist
\item
  Борелевская и лебегова \(\sigma\)-алгебры:
\end{enumerate}

\[
\mathcal B(\mathbb R^d)\subsetneq\mathcal L(\mathbb R^d)\subsetneq 2^{\mathbb R^d}.
\]

\subsection{8. Типичные дополнительные
вопросы}\label{ux442ux438ux43fux438ux447ux43dux44bux435-ux434ux43eux43fux43eux43bux43dux438ux442ux435ux43bux44cux43dux44bux435-ux432ux43eux43fux440ux43eux441ux44b}

\textbf{Вопрос. Что такое борелевское множество?}

Это элемент минимальной \(\sigma\)-алгебры, содержащей все открытые
множества:

\[
A\in\mathcal B(X).
\]

\textbf{Вопрос. Почему замкнутые множества борелевские?}

Потому что дополнение замкнутого множества открыто, а \(\sigma\)-алгебра
замкнута относительно дополнений.

\textbf{Вопрос. Почему в \(\mathbb R\) борелевская \(\sigma\)-алгебра
порождается лучами \((-\infty,x]\)?}

Потому что из лучей получаются полуинтервалы \((a,b]\), из них -
открытые интервалы, а открытые интервалы с рациональными концами
порождают обычную топологию прямой.

\textbf{Вопрос. Что значит, что мера полная?}

Это значит, что все подмножества измеримых множеств меры ноль тоже
измеримы.

\textbf{Вопрос. Чем борелевская мера Лебега отличается от меры Лебега?}

Борелевская мера Лебега задана только на \(\mathcal B(\mathbb R^d)\).
Полная мера Лебега задана на большей \(\sigma\)-алгебре
\(\mathcal L(\mathbb R^d)\), полученной пополнением.

\textbf{Вопрос. Может ли лебегово измеримое множество быть
неборелевским?}

Да. Например, некоторые подмножества канторова множества лебегово
измеримы, потому что лежат в множестве меры ноль, но не являются
борелевскими.

\textbf{Вопрос. Все ли подмножества \(\mathbb R\) становятся измеримыми
после пополнения?}

Нет. Существуют неизмеримые по Лебегу множества, например множества
Витали.

\subsection{9. Типичные
ошибки}\label{ux442ux438ux43fux438ux447ux43dux44bux435-ux43eux448ux438ux431ux43aux438}

\begin{enumerate}
\def\labelenumi{\arabic{enumi}.}
\tightlist
\item
  Считать борелевские и лебегово измеримые множества одним и тем же
  классом. Правильно:
\end{enumerate}

\[
\mathcal B(\mathbb R^d)\subsetneq\mathcal L(\mathbb R^d).
\]

\begin{enumerate}
\def\labelenumi{\arabic{enumi}.}
\setcounter{enumi}{1}
\item
  Не уточнять топологию в определении \(\mathcal B(X)\).
\item
  Думать, что пополнение добавляет все подмножества пространства. Оно
  добавляет только подмножества множеств меры ноль и их объединения со
  старыми измеримыми множествами.
\item
  Путать внешнюю меру и пополненную меру. Внешняя мера определена на
  всех подмножествах, но не является мерой на всех подмножествах.
\item
  Забывать, что борелевская мера Лебега не полна.
\item
  Думать, что всякое множество меры ноль счетно. Канторово множество
  несчетно и имеет меру ноль.
\item
  Использовать дополнение через верхний индекс. В этих конспектах:
\end{enumerate}

\[
\overline{A}=X\setminus A.
\]

\subsection{10. Связи с другими
билетами}\label{ux441ux432ux44fux437ux438-ux441-ux434ux440ux443ux433ux438ux43cux438-ux431ux438ux43bux435ux442ux430ux43cux438}

\textbf{Билет 01.} Борелевская \(\sigma\)-алгебра - важнейший пример
\(\sigma\)-алгебры.

\textbf{Билет 02.} Распределения случайных величин являются
вероятностными мерами на борелевских \(\sigma\)-алгебрах.

\textbf{Билет 03.} Борелевская мера Лебега строится как продолжение меры
с полуалгебры прямоугольников, а полная мера Лебега получается
пополнением.

\textbf{Билет 05.} Измеримость случайной величины означает, что
прообразы борелевских множеств измеримы.

\textbf{Билет 06.} Распределения случайных величин и векторов задаются
на \(\mathcal B(\mathbb R)\) и \(\mathcal B(\mathbb R^d)\).

\textbf{Билет 07.} Интеграл Лебега обычно строится по полной мере Лебега
или по абстрактной полной мере.

\subsection{11. Минимум на
удовлетворительно}\label{ux43cux438ux43dux438ux43cux443ux43c-ux43dux430-ux443ux434ux43eux432ux43bux435ux442ux432ux43eux440ux438ux442ux435ux43bux44cux43dux43e}

Нужно уметь сказать:

\begin{enumerate}
\def\labelenumi{\arabic{enumi}.}
\tightlist
\item
  Борелевская \(\sigma\)-алгебра:
\end{enumerate}

\[
\mathcal B(X)=\sigma(\text{открытые множества}).
\]

\begin{enumerate}
\def\labelenumi{\arabic{enumi}.}
\setcounter{enumi}{1}
\tightlist
\item
  Борелевская мера - мера на \(\mathcal B(X)\).
\item
  Борелевская мера Лебега на \(\mathbb R^d\) совпадает с объемом на
  прямоугольниках.
\item
  Полная мера содержит все подмножества множеств меры ноль.
\item
  Полная мера Лебега - пополнение борелевской меры Лебега.
\end{enumerate}

\subsection{12. Минимум на
хорошо}\label{ux43cux438ux43dux438ux43cux443ux43c-ux43dux430-ux445ux43eux440ux43eux448ux43e}

Кроме минимума на удовлетворительно, нужно:

\begin{enumerate}
\def\labelenumi{\arabic{enumi}.}
\tightlist
\item
  Доказать существование \(\mathcal B(X)\) как пересечения всех
  \(\sigma\)-алгебр, содержащих открытые множества.
\item
  Объяснить, почему замкнутые множества борелевские.
\item
  Назвать порождающие классы для \(\mathcal B(\mathbb R^d)\): шары,
  рациональные прямоугольники, полуоткрытые прямоугольники.
\item
  Сформулировать пополнение:
\end{enumerate}

\[
\overline{\mathcal F}^{\,\mu}
=
\{A\cup N:\ A\in\mathcal F,\ N\subset Z,\ \mu(Z)=0\}.
\]

\begin{enumerate}
\def\labelenumi{\arabic{enumi}.}
\setcounter{enumi}{4}
\tightlist
\item
  Объяснить различие между борелевскими и лебегово измеримыми
  множествами.
\end{enumerate}

\subsection{13. Что добавить для
отлично}\label{ux447ux442ux43e-ux434ux43eux431ux430ux432ux438ux442ux44c-ux434ux43bux44f-ux43eux442ux43bux438ux447ux43dux43e}

Для сильного ответа стоит добавить:

\begin{enumerate}
\def\labelenumi{\arabic{enumi}.}
\tightlist
\item
  Доказательство, что в пространстве со счетной базой борелевская
  \(\sigma\)-алгебра порождена базой.
\item
  Доказательство, что \(\mathcal B(\mathbb R)=\sigma\{(-\infty,x]\}\).
\item
  Эквивалентное описание пополнения через симметрическую разность:
\end{enumerate}

\[
E\triangle A\subset Z,\qquad \mu(Z)=0.
\]

\begin{enumerate}
\def\labelenumi{\arabic{enumi}.}
\setcounter{enumi}{3}
\tightlist
\item
  Обоснование, почему борелевская мера Лебега не полна, через
  подмножества канторова множества.
\item
  Цепочку:
\end{enumerate}

\[
\mathcal B(\mathbb R^d)\subsetneq\mathcal L(\mathbb R^d)\subsetneq 2^{\mathbb R^d}.
\]

\subsection{14. Устная версия ответа на 5
минут}\label{ux443ux441ux442ux43dux430ux44f-ux432ux435ux440ux441ux438ux44f-ux43eux442ux432ux435ux442ux430-ux43dux430-5-ux43cux438ux43dux443ux442}

Борелевская \(\sigma\)-алгебра строится из топологии. Если \(X\) -
топологическое пространство, то

\[
\mathcal B(X)=\sigma\{G\subset X:\ G\text{ открыто}\}.
\]

Это минимальная \(\sigma\)-алгебра, содержащая все открытые множества.
Ее элементы называются борелевскими множествами. Замкнутые множества
тоже борелевские, потому что они являются дополнениями открытых
множеств.

В \(\mathbb R^d\) борелевскую \(\sigma\)-алгебру можно порождать разными
классами: открытыми шарами, замкнутыми множествами, прямоугольниками,
полуоткрытыми прямоугольниками. Если взять прямоугольники с
рациональными координатами, они образуют счетную базу топологии, поэтому
тоже порождают \(\mathcal B(\mathbb R^d)\).

Борелевская мера - это мера на \(\mathcal B(X)\). Борелевская мера
Лебега на \(\mathbb R^d\) - это мера, которая на полуоткрытых
прямоугольниках равна объему:

\[
\lambda_d^B\left(\prod_{j=1}^d[a_j,b_j)\right)
=
\prod_{j=1}^d(b_j-a_j).
\]

Но борелевская мера Лебега не полна: могут быть подмножества борелевских
множеств меры ноль, которые сами не являются борелевскими.

Поэтому вводят пополнение меры. Мера \(\mu\) на \((\Omega,\mathcal F)\)
называется полной, если

\[
N\subset Z,\quad Z\in\mathcal F,\quad \mu(Z)=0
\Rightarrow N\in\mathcal F.
\]

Пополнение добавляет к \(\mathcal F\) все множества вида

\[
A\cup N,
\]

где \(A\in\mathcal F\), а \(N\) лежит в измеримом множестве меры ноль.
Формально:

\[
\overline{\mathcal F}^{\,\mu}
=
\{A\cup N:\ A\in\mathcal F,\ N\subset Z,\ Z\in\mathcal F,\ \mu(Z)=0\}.
\]

Полная мера Лебега - это пополнение борелевской меры Лебега. Поэтому

\[
\mathcal B(\mathbb R^d)\subsetneq\mathcal L(\mathbb R^d)\subsetneq 2^{\mathbb R^d}.
\]

Первое включение строго, потому что есть лебегово измеримые
неборелевские множества, например некоторые подмножества канторова
множества. Второе включение строго, потому что существуют неизмеримые
множества, например множества Витали.

\subsection{15. Если осталось 2 минуты перед
экзаменом}\label{ux435ux441ux43bux438-ux43eux441ux442ux430ux43bux43eux441ux44c-2-ux43cux438ux43dux443ux442ux44b-ux43fux435ux440ux435ux434-ux44dux43aux437ux430ux43cux435ux43dux43eux43c}

Запомнить:

\[
\mathcal B(X)=\sigma(\text{открытые множества}).
\]

Борелевская мера - мера на \(\mathcal B(X)\).

Борелевская мера Лебега:

\[
\lambda_d^B\left(\prod_{j=1}^d[a_j,b_j)\right)
=
\prod_{j=1}^d(b_j-a_j).
\]

Полная мера:

\[
N\subset Z,\quad \mu(Z)=0
\Rightarrow
N\text{ измеримо}.
\]

Пополнение:

\[
\overline{\mathcal F}^{\,\mu}
=
\{A\cup N:\ A\in\mathcal F,\ N\subset Z,\ Z\in\mathcal F,\ \mu(Z)=0\}.
\]

Главная мысль:

\[
\mathcal B(\mathbb R^d)\subsetneq\mathcal L(\mathbb R^d)\subsetneq 2^{\mathbb R^d}.
\]

Борелевские множества - не то же самое, что все лебегово измеримые
множества.

\newpage

\end{document}
